\section{Beyond the Torah: A Library Grows Book by Book}

After the Pentateuch, further Hebrew and Aramaic books entered Greek through independent projects. Dates and places are often recoverable only in broad ranges. Many translations probably arose in Egypt during the second and first centuries \textsc{bc}; some evidence points elsewhere, and revisions could occur far from the place of first translation. “The Septuagint translators” is therefore least accurate precisely where the corpus becomes largest.

\subsection{The anchor supplied by Ben Sira}

The Greek prologue to Sirach is unusually valuable because a translator speaks in his own voice. He identifies the Hebrew author as his grandfather, reports arriving in Egypt in the thirty-eighth year of King Euergetes---conventionally understood as 132 \textsc{bc}---and explains the difficulty of rendering Hebrew into another language. He also refers to the Law, the Prophets, and other ancestral books in Greek.

This is strong evidence that several classes of Jewish Scripture had Greek forms by the late second century \textsc{bc}. It is not a table of contents, a declaration that every later canonical book had been translated, or proof of a closed tripartite canon. The repeated, flexible description identifies an inherited scriptural world without listing its borders.

\subsection{Translations, compositions, and additions}

Not everything transmitted within later Septuagint manuscripts is a translation from Hebrew. The wider Greek corpus contains at least four literary relationships:

\begin{threestable}{Relationship}{Representative cases}{Historical consequence}
Greek translation of a Semitic book & Pentateuch, many Prophets, Psalms, and other books. & Requires study of translation technique and a recoverable Hebrew or Aramaic source text.\\
Greek form of a different literary edition & Jeremiah and parts of Samuel--Kings are central test cases. & Difference may precede translation and illuminate the book's literary growth.\\
Greek expansions or alternative forms & Esther and Daniel; the longer and shorter Greek forms of Tobit illustrate further plurality. & “Addition” and “original” must be defined by book and witness; later codices may transmit hybrid texts.\\
Work composed in Greek & Wisdom of Solomon and 2 Maccabees are leading examples. & Calling the whole collection a “translation” erases Greek Jewish authorship.\\
\end{threestable}

Modern Catholic “deuterocanonical” and Protestant “Apocrypha” classifications concern later canon reception. They cannot simply be projected onto Jewish authors or the first Greek scrolls. Alison Salvesen's study of these books emphasizes both their Jewish origins and the variety of contents and groupings in later Christian codices. Some works once existed in Hebrew or Aramaic and survive substantially through Greek and daughter versions; others were Greek compositions from the outset. Their preservation by Christians proves reception, not one ancient Alexandrian list.

\subsection{A chronology with different clocks}

Books could be translated, revised, expanded, and collected at different times. A Greek version's latest datable vocabulary or historical allusion may help place it, but translation Greek is conservative and source-bound. A manuscript's date establishes when that copy was made, not when its underlying translation originated. A patristic citation establishes use by that author, not the first appearance of the cited form.

The result is a layered chronology:

\begin{itemize}
\item third century \textsc{bc}: the Greek Pentateuch supplies the foundational model and vocabulary;
\item second and first centuries \textsc{bc}: many further translations and Greek Jewish compositions appear, with Sirach's prologue providing a firm late-second-century horizon for substantial prior activity;
\item first centuries \textsc{bc} and \textsc{ad}: existing versions are revised toward emerging Hebrew reference forms, while Jewish authors continue to cite and interpret Greek Scripture;
\item second and third centuries \textsc{ad}: new Jewish Greek translations and Christian editorial work multiply the recoverable forms;
\item fourth and fifth centuries: large Christian codices gather texts whose histories began in separate scrolls.
\end{itemize}

No point on this timeline is “the completion of the Septuagint” in the way a modern publishing house completes one edition. The absence of a single finishing date is one of the history's principal conclusions.

\claimcheck{An “Alexandrian canon”}{Neither Aristeas's Law, Ben Sira's threefold description, nor the great Christian codices supplies a fixed Alexandrian Jewish table of contents. The codices differ among themselves and belong to a later Christian setting. Canon history must be argued from explicit judgments and patterns of use, not inferred from the modern contents page of a printed Septuagint.}
